\documentclass[11pt]{article}
\usepackage[a4paper,margin=0.65in]{geometry}
\usepackage[T1]{fontenc}
\usepackage{amsmath,amssymb,amsfonts}
\usepackage{graphicx}
\usepackage{pdfpages}
\usepackage[english,shorthands=off]{babel}
\usepackage{fancyhdr}
\usepackage[bookmarks=false,colorlinks=false]{hyperref}
\usepackage[protrusion=true,expansion=false]{microtype}
\usepackage{array}
\usepackage{longtable}
\usepackage{multirow}
\emergencystretch=3em
\tolerance=600
\providecommand{\modd}{\mathrm{mod}}
\providecommand{\Disc}{\mathrm{Disc}}
\providecommand{\canont}{}
\providecommand{\asterism}{*}
\providecommand{\Hessian}{H}
\providecommand{\tome}{}
\providecommand{\page}{p.}
\pagestyle{fancy}\fancyhf{}\fancyhead[L]{Pages 451--475}\fancyhead[R]{Cayley --- Collected Papers, Vol. XIII}\renewcommand{\headrulewidth}{0.4pt}
\setlength{\headheight}{15pt}

\newcommand{\paper}[3]{%
  \clearpage
  \begin{center}
    {\large\bfseries #2}\\[0.5em]
    {\small [#1]}\\[0.3em]
    {\small #3}
  \end{center}
  \vspace{0.5em}
}

\begin{document}

\setcounter{page}{425}

% ====================================================================
% PAPER 947 (PDF pp. 451-455, printed pp. 425-429)
% ON A SYSTEM OF TWO TETRADS OF CIRCLES; AND OTHER SYSTEMS OF TWO TETRADS.
% ====================================================================
\paper{947}{ON A SYSTEM OF TWO TETRADS OF CIRCLES; AND OTHER SYSTEMS OF TWO TETRADS.}{[From the \textit{Proceedings of the Cambridge Philosophical Society}, vol.\ \textsc{viii}.\ (1893), pp.\ 54--59.]}

% ----- printed p. 425 (PDF 451) -----
\noindent The investigations of the present paper were suggested to me by Mr Orr's paper, ``The Contacts of certain Systems of Circles,'' \textit{Proc.\ Camb.\ Phil.\ Soc.}, vol.\ \textsc{vii}.

\medskip
\noindent\textbf{1.} It is possible to find \textit{in plano} two tetrads of circles, or say four red circles and four blue circles, such that each red circle touches each blue circle: in fact, counting the constants, a circle depends upon 3 constants, or say it has a capacity $= 3$; the capacity of the eight circles is thus $= 24$; and the postulation or number of conditions to be satisfied is $= 16$: the resulting capacity of the system is thus \textit{prim\^a facie}, $16 - 24 = 8$. It will, however, appear that, in the system considered, the true value is $= 9$.

\medskip
\noindent\textbf{2.} The \textit{prim\^a facie} value of the capacity being $= 8$, we are not at liberty to assume at pleasure three circles of the system. And, in fact, assuming at pleasure say 3 red circles, then touching each of these we have 8 circles, forming $\tfrac{1}{24}\,8 \cdot 7 \cdot 6 \cdot 5$, $= 70$, tetrads of circles: taking at random any one of these tetrads for the blue circles, the remaining red circle has to be determined so as to touch each of the four blue circles, that is, by four instead of three conditions; and there is not in general any red circle satisfying these four conditions. But the 8 tangent circles do not stand to each other in a relation of symmetry, but form in fact four pairs of circles; and it is possible out of the 70 tetrads to select (and that in 6 ways) a tetrad of blue circles, such that there exists a fourth red circle touching each of these four blue circles. We have thus a system depending upon 3 arbitrary circles, and for which, therefore, the capacity is $= 9$. It is (as is known) possible, in quite a

% ----- printed p. 426 (PDF 452) -----
different manner, out of the 70 tetrads to select (and that in 8 ways) a tetrad of blue circles such that there exists a fourth red circle touching each of these four blue circles---but the present paper relates exclusively to the first-mentioned 6 tetrads and not to these 8 tetrads.

\medskip
\noindent\textbf{3.} I consider, in the first instance, a particular case in which the three red circles are not all of them arbitrary, but have a capacity $9 - 1$, $= 8$; and pass from this to the general case where the capacity is $= 9$. Calling the red circles 1, 2, 3 and 4; I start with the circles 1, and 2 arbitrary, and 3 a circle equal to 2: the radical axis, or common chord, of the circles 2 and 3 is thus a line bisecting at right angles the line joining the centres of the circles 2 and 3, say this is the line $\Omega$. We have then four circles, each having its centre in the line $\Omega$ and touching the circles 1 and 2: in fact, the locus of the centre of a circle touching the circles 1 and 2 is a pair of conics, each of them having for foci the centres of these circles: the line $\Omega$ meets each of these conics in two points, and there are thus on the line $\Omega$ four points, each of them the centre of a circle touching the circles 1 and 2. But the equal circles 2 and 3 are symmetrically situate in regard to the line $\Omega$; and it is obvious that the four circles, having their centres on the line $\Omega$, will each of them also touch the circle 3; we have thus the four blue circles, each of them with its centre on the line $\Omega$ and touching each of the red circles 1, 2 and 3. And it is moreover clear that, taking the red circle 4 equal to 1 and situate symmetrically therewith in regard to the line $\Omega$, then this circle 4 will touch each of the blue circles: so that we have here the four blue circles, each of them touching the four red circles. As already mentioned, the blue circles have their centre on the line $\Omega$, that is, the line $\Omega$ is a common orthotomic of the four blue circles.

\medskip
\noindent\textbf{4.} By inverting in regard to an arbitrary circle we pass to the general case; the line $\Omega$ becomes thus a circle $\Omega$, orthotomic to each of the blue circles.

Starting \textit{ab initio}, we have here at pleasure the red circles 1, 2, 3: the circle $\Omega$ is a circle having for centre a centre of symmetry of the circles 2 and 3, and passing through the points of intersection (real or imaginary) of these two circles; the circles 2 and 3 are thus the inverses (or say the images) each of the other in regard to the circle $\Omega$. We can then find 4 circles each of them orthotomic to $\Omega$, and touching the circles 1 and 2: but a circle orthotomic to $\Omega$ is its own inverse or image in regard to $\Omega$; and it will thus touch the circle 3 which is the image of 2 in regard to $\Omega$. We have thus the four blue circles each of them touching the red circles 1, 2 and 3; and then, taking the red circle 4 as the inverse or image of 1 in regard to $\Omega$, this circle 4 will also touch each of the blue circles. Thus starting with the arbitrary red circles 1, 2, 3, we find the four blue circles and the remaining red circle 4, such that each of the blue circles touches each of the red circles. Since in the construction we group together at pleasure the two circles 2, 3 (out of the three circles 1, 2, 3) and use at pleasure either of the two centres of symmetry, it appears that the number of ways in which the figure might have been completed is $= 6$.

% ----- printed p. 427 (PDF 453) -----
\medskip
\noindent\textbf{5.} The blue circles have a common orthotomic circle $\Omega$, that is, the radical axis or common chord of each two of the blue circles passes through one and the same point, the centre of the circle $\Omega$. The figure is symmetrical in regard to the red and blue circles respectively, and thus the red circles have a common orthotomic circle $\Omega'$, that is, the radical axis or common chord of each two of the red circles passes through one and the same point, the centre of the circle $\Omega'$.

\medskip
\noindent\textbf{6.} Projecting stereographically on a spherical surface, the four red circles and the four blue circles become circles of the sphere; and then making the general homographic transformation, they become plane sections of a quadric surface; we have thus the theorem that on a given quadric surface it is possible to find four red sections and four blue sections such that each blue section touches each red section; and moreover the capacity of the system is $= 9$; viz.\ 3 of the red sections may be assumed at pleasure. But (as is well known) the theory of the tangency of plane sections of a quadric surface is far more simple than that of the tangency of circles: the condition in order that two sections may touch each other is simply the condition that the line of intersection of the two planes shall touch the quadric surface. And we construct, as follows, the sections touching each of three given sections: say the given sections are 1, 2, 3; through the sections 1 and 2 we have two quadric cones having for vertices say the points $d_{12}$ and $i_{12}$ (direct and inverse centres of the two sections): similarly, through the sections 1 and 3 we have two quadric cones vertices $d_{13}$ and $i_{13}$ respectively, and through the sections 2 and 3 we have two quadric cones vertices $d_{23}$ and $i_{23}$ respectively; the points $d_{12}, i_{12}, d_{13}, i_{13}, d_{23}, i_{23}$, lie three and three in four intersecting lines or axes, viz.\ these are $d_{23}d_{31}d_{12}$, $d_{23}i_{31}i_{12}$, $d_{31}i_{12}i_{23}$, $d_{12}i_{23}i_{31}$ respectively. Through any one of these axes, say $d_{23}d_{31}d_{12}$, we may draw to the quadric surface two tangent planes each touching the three cones which have their vertices in the points $d_{23}, d_{31}, d_{12}$ respectively; and the section by either tangent plane is thus a section touching each of the three given sections 1, 2, 3; we have thus the eight tangent sections of these three sections.

\medskip
\noindent\textbf{7.} Taking as three of the red sections the arbitrary sections 1, 2, 3; and grouping together two at pleasure of these sections, say 2 and 3; we may take for the blue sections the two sections through the axis $d_{23}d_{31}d_{12}$, and those through the axis $d_{23}i_{31}i_{12}$; we have thus the four blue sections touching each of the given red sections 1, 2, 3; and this being so, there exists a remaining red section 4 touching each of the blue sections; we have thus the four blue sections touching each of the red sections 1, 2, 3 and 4. This implies that the vertices or points $d_{24}$ and $d_{34}$ lie on the axis $d_{23}d_{31}d_{12}$, and that the vertices or points $i_{24}$ and $i_{34}$ lie on the axis $d_{23}i_{31}i_{12}$; or, what is the same thing, that the four sections 1, 2, 3, 4 have in common an axis $d_{23}d_{31}d_{12}d_{24}d_{34}$ and also an axis $d_{23}i_{31}i_{12}i_{24}i_{34}$.

\medskip
\noindent\textbf{8.} If the quadric surface be a flat surface (\textit{surface aplatie}) or conic, then the red sections become chords of the conic; the axes are lines in the plane of the conic, and thus the tangent planes through an axis each coincide with the plane of the conic, and it would at first sight appear that any theorem as to tangency becomes nugatory. But this is not so; comparing with the last preceding paragraph, we still have the theorem: on a given conic, taking at pleasure any three chords

% ----- printed p. 428 (PDF 454) -----
1, 2, 3, it is possible to find a fourth chord 4, such that the four chords have in common an axis $d_{23}d_{31}d_{12}d_{24}d_{34}$ and also an axis $d_{23}i_{31}i_{12}i_{24}i_{34}$. And the analytical theory (although somewhat complex) is extremely interesting. Considering the conic $xz - y^2 = 0$, the coordinates of a point on the conic are given by $x : y : z = 1 : \theta : \theta^2$, or say any point of the conic is determined by its parameter $\theta$; and this being so, considering any three chords 1, 2, 3, I take for the two extremities of 1 the values $\epsilon, \zeta$; for those of 2 the values $\alpha, \beta$; and for those of 3 the values $\gamma, \delta$; the remaining chord 4 is to be determined as above, and I take for its two extremities the values $E, Z$.

\medskip
\noindent\textbf{9.} Starting with the chords 1, 2, 3, we have each of the points $d_{23}$, \&c., as the intersection of two lines, viz.\ these are
\[
\begin{aligned}
&\text{for } d_{23}\colon\quad
\begin{cases} \alpha\beta\,z - y\,(\alpha + \beta) + x = 0, \\ \gamma\delta\,z - y\,(\gamma + \delta) + x = 0; \end{cases}
\qquad
\text{for } d_{31}\colon\quad
\begin{cases} \alpha\gamma\,z - y\,(\alpha + \gamma) + x = 0, \\ \beta\delta\,z - y\,(\beta + \delta) + x = 0; \end{cases} \\[4pt]
&\text{for } d_{12}\colon\quad
\begin{cases} \alpha\epsilon\,z - y\,(\alpha + \epsilon) + x = 0, \\ \beta\zeta\,z - y\,(\beta + \zeta) + x = 0; \end{cases}
\qquad
\text{for } i_{12}\colon\quad
\begin{cases} \alpha\zeta\,z - y\,(\alpha + \zeta) + x = 0, \\ \beta\epsilon\,z - y\,(\beta + \epsilon) + x = 0; \end{cases}
\end{aligned}
\]
and similarly for $i_{23}$, $i_{31}$. We thence find without difficulty for the axis $d_{23}d_{31}d_{12}$ the equation
\[
\resizebox{\textwidth}{!}{$\displaystyle
x\,\bigl\{\alpha\beta\,(\delta\epsilon - \gamma\zeta) + \gamma\delta\,(\zeta\alpha - \epsilon\beta) + \epsilon\zeta\,(\beta\gamma - \alpha\delta)\bigr\}
+ y\,\bigl\{(\beta - \alpha)(\gamma\delta - \epsilon\zeta) + (\delta - \gamma)(\epsilon\zeta - \alpha\beta) + (\zeta - \epsilon)(\alpha\beta - \gamma\delta)\bigr\}
+ z\,\bigl\{-(\delta\epsilon - \gamma\zeta) - (\zeta\alpha - \epsilon\beta) - (\beta\gamma - \alpha\delta)\bigr\} = 0;
$}
\]
the equation of the axis $d_{23}i_{31}i_{12}$ is obtained herefrom by the interchange of $\epsilon$ and $\zeta$.

\medskip
\noindent\textbf{10.} The points $d_{24}$ and $d_{34}$ will lie upon the first-mentioned axis if only $d_{34}$ lies upon this axis, viz.\ if we have
\[
\resizebox{\textwidth}{!}{$\displaystyle
\bigl[\alpha\beta\,(\delta\epsilon - \gamma\zeta) + \gamma\delta\,(\zeta\alpha - \epsilon\beta) + \epsilon\zeta\,(\beta\gamma - \alpha\delta)\bigr]\,(\beta + E - \alpha - Z)
+ \bigl[(\beta - \alpha)(\gamma\delta - \epsilon\zeta) + (\delta - \gamma)(\epsilon\zeta - \alpha\beta) + (\zeta - \epsilon)(\alpha\beta - \gamma\delta)\bigr]\,(\beta E - \alpha Z)
+ \bigl[-(\delta\epsilon - \gamma\zeta) - (\zeta\alpha - \epsilon\beta) - (\beta\gamma - \alpha\delta)\bigr]\,\bigl[\,\cdots\,\bigr] = 0.
$}
\]
Reducing this equation, the factor $\alpha - \beta$ divides out, and we finally obtain
\[
A + BE + CZ + DEZ = 0;
\]
say this is
\[
A + BE + CZ + DEZ = 0,
\]
where
\[
\begin{aligned}
A &= \,\cdot\, , \\
B &= -(\gamma - \alpha)\,\beta\delta\,. \quad +\; (\alpha\beta - \gamma\delta)\,\zeta \quad +\; (\beta - \delta)\,\epsilon\zeta, \\
C &= -(\beta - \delta)\,\alpha\gamma + (\alpha\beta - \gamma\delta)\,\epsilon \quad +\; (\gamma - \alpha)\,\epsilon\zeta, \\
D &= -(\alpha\delta - \beta\gamma) - (\beta - \delta)\,\epsilon - (\gamma - \alpha)\,\zeta,
\end{aligned}
\]
viz.\ this is the condition for the existence of the axis $d_{23}d_{31}d_{12}d_{24}d_{34}$.

% ----- printed p. 429 (PDF 455) -----
We interchange herein $\epsilon$, $\zeta$ and also $E$, $Z$, and we thus obtain
\[
A' + C'E + B'Z + D'EZ = 0,
\]
where
\[
\begin{aligned}
A' &= \,\cdot\, , \\
B' &= \,\cdot\, +\, (\beta - \delta)\,\alpha\gamma\,\zeta + (\gamma - \alpha)\,\beta\delta\,\epsilon + (\alpha\beta - \gamma\delta)\,\epsilon\zeta, \\
C' &= (\alpha\beta - \gamma\delta)\,\epsilon \quad +\; (\zeta - \delta)\,\cdot\, , \\
D' &= \,\cdot\, +\, (\alpha\beta - \gamma\delta)\,\zeta + (\gamma - \alpha)\,\cdot\, ,
\end{aligned}
\]
viz.\ this is the condition for the existence of the axis $d_{23}i_{31}i_{12}i_{24}i_{34}$.

\medskip
\noindent\textbf{11.} I remark that we have
\[
A' = A + \mathrm{a}\,(\epsilon - \zeta), \quad B' = B + \mathrm{b}\,(\epsilon - \zeta),
\]
\[
C' = C + \mathrm{c}\,(\epsilon - \zeta), \quad D' = D + \mathrm{d}\,(\epsilon - \zeta),
\]
where
\[
\begin{aligned}
\mathrm{a} &= (\beta - \delta)\,\alpha\gamma - (\gamma - \alpha)\,\beta\delta = \alpha\beta\,(\gamma + \delta) - \gamma\delta\,(\alpha + \beta), \\
\mathrm{b} &= \mathrm{c} = \gamma\delta - \alpha\beta, \\
\mathrm{d} &= \alpha + \beta - \gamma - \delta,
\end{aligned}
\]
and further that
\[
B - C = \alpha\beta\,(\gamma + \delta) - \gamma\delta\,(\alpha + \beta) + (\gamma\delta - \alpha\beta)\,(\epsilon + \zeta) + (\alpha + \beta - \gamma - \delta)\,\epsilon\zeta;
\]
say this is $\Pi$.

\medskip
\noindent\textbf{12.} It thus appears that, for the determination of $E$, $Z$, we have
\[
A + BE + CZ + DEZ = 0,
\]
\[
A' + C'E + B'Z + D'EZ = 0.
\]
Eliminating $Z$, we find
\[
\frac{A + BE}{A' + C'E} = \frac{C + DE}{B' + D'E},
\]
that is,
\[
(AC' - A'C) + (AD' - A'D + BB' - CC')\,E + (BD' - B'D)\,E^2 = 0;
\]
upon reducing the coefficients of this equation it appears that they contain each of them the factor $\Pi$, and throwing out this factor, the equation is
\[
\epsilon\,\bigl[\alpha\beta\,(\gamma - \delta) + \gamma\delta\,(\beta - \alpha) + (\alpha\delta - \beta\gamma)\,\zeta\bigr] + \cdots = 0;
\]
this contains obviously the factor $E - \epsilon$, or throwing out this factor, we have for $E$ the simple equation
\[
\bigl[\alpha\beta\,(\gamma - \delta) + \gamma\delta\,(\beta - \alpha)\bigr] + (\alpha\delta - \beta\gamma)\,(E + \zeta) + (\beta - \alpha + \gamma - \delta)\,\zeta E = 0.
\]
In a similar manner it may be shown that the two equations give for $Z$ the like simple equation
\[
\bigl[\alpha\beta\,(\gamma - \delta) + \gamma\delta\,(\beta - \alpha)\bigr] + (\alpha\delta - \beta\gamma)\,(Z + \epsilon) + (\beta - \alpha + \gamma - \delta)\,\epsilon Z = 0,
\]
viz.\ starting from the chords 1, 2, 3 which depend on the parameters $(\epsilon, \zeta)$, $(\alpha, \beta)$, $(\gamma, \delta)$ respectively, these last two equations give the parameters $(E, Z)$ of the chord 4.

% ====================================================================
% PAPER 948 (PDF pp. 456-475, printed pp. 430-449)
% REPORT OF A COMMITTEE APPOINTED FOR THE PURPOSE OF CARRYING ON THE
% TABLES CONNECTED WITH THE PELLIAN EQUATION FROM THE POINT WHERE THE
% WORK WAS LEFT BY DEGEN IN 1817.
% ====================================================================
\paper{948}{REPORT OF A COMMITTEE APPOINTED FOR THE PURPOSE OF CARRYING ON THE TABLES CONNECTED WITH THE PELLIAN EQUATION FROM THE POINT WHERE THE WORK WAS LEFT BY DEGEN IN 1817.}{[From the \textit{British Association Report}, (1893), pp.\ 73--120.]}

% ----- printed p. 430 (PDF 456) -----
\noindent We have, on the Pellian Equation, Degen's tables, the title of which is ``Canon Pellianus sive Tabula simplicissimam {\ae}quationis celebratissim{\ae} $y^2 = ax^2 + 1$ solutionem pro singulis numeri dati valoribus ab 1 usque ad 1000 in numeris rationalibus iisdemque integris exhibens.'' Autore Carolo Ferdinando Degen. Hafni{\ae}, apud Gerhardum Bonnierum, MDCCCXVII., 8vo. Introductio, pp.\ v---xxiv. Tabula I.\ Solutionem {\ae}quationis $y^2 - ax^2 - 1 = 0$ exhibens, pp.\ 3---106. Tabula II.\ Solutionem {\ae}quationis $y^2 - ax^2 + 1 = 0$, quotiescunque valor ipsius $a$ talem admiserit, exhibens, pp.\ 109---112.

The mode of calculation is explained in the Introduction, and illustrated by the examples of the numbers 209, 173.

As to the first of these, the entry in Table I.\ is
\begin{center}
\begin{tabular}{c|l}
\multirow{4}{*}{$209$} & $14,\ 2,\ 5,\ 3,\ (2)$ \\
 & $1,\ 13,\ 5,\ 8,\ 11$ \\
 & $3220$ \\
 & $46551$ \\
\end{tabular}
\end{center}
where the first line gives the expression of $\sqrt{209}$ as a continued fraction, viz.\ we have
\[
\sqrt{209} = 14 + \cfrac{1}{2 +} \cfrac{1}{5 +} \cfrac{1}{3 +} \cfrac{1}{2 +} \cfrac{1}{3 +} \cfrac{1}{5 +} \cfrac{1}{2 +} \cfrac{1}{28 +} \cfrac{1}{2 +}\;\&\text{c.}
\]

% ----- printed p. 431 (PDF 457) -----
the denominators being 2, 5, 3, (2), 3, 5, 2, then 28, which is the double of the integer part 14, and then again 2, 5, 3, (2), 3, 5, 2, and so on, the parentheses of the (2) being used to indicate that this is the middle term of the period.

The second row gives auxiliary numbers occurring in the calculation of the first row and having a meaning, as will presently appear. Observe that the 11 which comes under the (2) should also be printed in parentheses (11), but this is not done.

The process for the calculation of the $x$, $y$ is as follows:
\begin{center}
\begin{tabular}{|c|c|c|c|}
\hline
\multicolumn{4}{|c|}{$209$} \\
\hline
 & $1$ & $0$ & $+ 1$ \\
$14$ & $14$ & $1$ & $-13$ \\
$2$ & $29$ & $2$ & $+ 5$ \\
$5$ & $159$ & $11$ & $- 8$ \\
$3$ & $506$ & $35$ & $+ (11)$ \\
$(2)$ & $1171$ & $81$ & $-13$ \\
$3$ & $4019$ & $278$ & $+ 5$ \\
$5$ & $21266$ & $1471$ & $- 8$ \\
$2$ & & $3220$ & \\
$28$ & $46551$ & & $+ 1$ \\
\hline
\end{tabular}
\end{center}
viz.\ writing down as a first column the numbers of the first row, and beginning the second column with 1, 14 (14 the number at the head of the first column), and the third column with 0, 1, we calculate the numbers of the second column, $29 = 2 \cdot 14 + 1$, $159 = 5 \cdot 29 + 14$, $506 = 3 \cdot 159 + 29$, \&c., and the numbers of the third column in like manner, $2 = 2 \cdot 1 + 0$, $11 = 5 \cdot 2 + 1$, $35 = 3 \cdot 11 + 2$, \&c.; and then writing down as a fourth column the numbers of the second row with the signs $+$, $-$ alternately, we have a series of equations $y^2 - ax^2 = \pm A$, viz.
\[
1^2 - 209 \cdot 0^2 = + 1,
\]
\[
14^2 - 209 \cdot 1^2 = - 13,
\]
\[
29^2 - 209 \cdot 2^2 = + 5,
\]
the last of them being
\[
(46551)^2 - 209\,(3220)^2 = + 1,
\]
this last corresponding as above to the value $+1$, and the numbers 46551 and 3220 being accordingly the $y$ and $x$ given in the fourth and third rows of the table.

% ----- printed p. 432 (PDF 458) -----
As to the second of the foregoing numbers, 173, the only difference is that the period has a double middle term, viz.\ the entry in the Table I.\ is
\begin{center}
\begin{tabular}{c|l}
\multirow{4}{*}{$173$} & $13,\ 6,\ (1,\ 1)$ \\
 & $1,\ 4,\ (13,\ 13)$ \\
 & $190060$ \\
 & $2499849$ \\
\end{tabular}
\end{center}

The first row gives the expression of $\sqrt{173}$, viz.\ that is
\[
\sqrt{173} = 13 + \cfrac{1}{6 +} \cfrac{1}{1 +} \cfrac{1}{1 +} \cfrac{1}{6 +} \cfrac{1}{26 +}\;\&\text{c.}
\]
the denominators being 6, 1, 1, 6, then 26 (the double of the integer part 13), and then again 6, 1, 1, 6, and so on. In the second row I remark that Degen prints the parentheses (13, 13) for the double middle term.

The process for the calculation of the $x$, $y$ is similar to that in the former case, viz.\ we have
\begin{center}
\begin{tabular}{|c|c|c|c|}
\hline
\multicolumn{4}{|c|}{$173$} \\
\hline
 & $1$ & $0$ & $+ 1$ \\
$13$ & $13$ & $1$ & $- 4$ \\
$6$ & $79$ & $6$ & $+ 13$ \\
$(1,\,1)$ & $92$ & $7$ & $- 1$ \\
 & $171$ & $13$ & $+ 13$ \\
$6$ & & & \\
$26$ & $1118$ & $85$ & $+ 4$ \\
\hline
\end{tabular}
\end{center}
where the second and third columns begin 1, 13 and 0, 1 respectively, and the remaining terms are calculated $79 = 6 \cdot 13 + 1$, $92 = 1 \cdot 79 + 13$, \&c., and $6 = 6 \cdot 1 + 0$, $7 = 1 \cdot 6 + 1$, \&c.; and then writing down as a fourth column the terms of the second row with the signs $+$, $-$ alternately, we have
\[
1^2 - 173 \cdot 0^2 = + 1,
\]
\[
13^2 - 173 \cdot 1^2 = - 4,
\]
\[
79^2 - 173 \cdot 6^2 = + 13,
\]
the last equation being
\[
(1118)^2 - 173\,(85)^2 = - 1,
\]

% ----- printed p. 433 (PDF 459) -----
the term for the last equation being always in a case such as the present one, not $+ 1$, but $- 1$. The final numbers 1118, 85 are consequently entered not in Table I., but in Table II., viz.\ the entry in this table is
\begin{center}
\begin{tabular}{c|c}
\multirow{2}{*}{$173$} & $85$ \\
 & $1118$ \\
\end{tabular}
\end{center}
and thence we calculate the numbers $y$, $x$ of Table I., viz.\ these are
\[
2499849 = 2 \cdot (1118)^2 + 1,
\]
\[
190060 = 2 \cdot 1118 \cdot 85.
\]

Generally Table II.\ gives for each value of $a$, comprised therein, values of $x$, $y$, such that $y^2 = ax^2 - 1$, and then writing $y_1 = 2y^2 + 1$, $x_1 = 2xy$, we have
\[
y_1^2 = (2ax^2 - 1)^2 = 4a^2 x^4 - 4ax^2 + 1 = a \cdot 4x^2\,(ax^2 - 1) + 1 = a x_1^2 + 1,
\]
so that $x_1$, $y_1$ are for the same value of $a$ the values of $x$, $y$ in Table I.

It is to be remarked that the heading of Table II.\ is not perfectly accurate, for it purports to give for every value of $a$, for which a solution exists, a solution of the equation $y^2 = ax^2 - 1$. What it really gives is the solution for each value of $a$ for which the period has a double middle term. But if $a = \alpha^2 + 1$, then obviously we have a solution $y = \alpha$, $x = 1$, and for any such value of $a$ the period has a single middle term, viz.\ the entry in Table I.\ is
\begin{center}
\begin{tabular}{c|l}
\multirow{4}{*}{$\alpha^2 + 1$} & $\alpha,\ (2\alpha)$ \\
 & $1,\ 1$ \\
 & $2\alpha$ \\
 & $2\alpha^2 + 1$ \\
\end{tabular}
\end{center}
and we, in fact, have
\begin{center}
\begin{tabular}{|c|c|c|c|}
\hline
\multicolumn{4}{|c|}{$\alpha^2 + 1$} \\
\hline
 & $1$ & $0$ & $- 1$ \\
$\alpha$ & $\alpha$ & $1$ & $+ 1$ \\
$(2\alpha)$ & $2\alpha^2 + 1$ & $2\alpha$ & $+ 1$ \\
\hline
\end{tabular}
\end{center}
that is,
\[
1^2 - (\alpha^2 + 1) \cdot 0^2 = + 1,
\]
\[
\alpha^2 - (\alpha^2 + 1) \cdot 1^2 = - 1,
\]
\[
(2\alpha^2 + 1)^2 - (\alpha^2 + 1)\,(2\alpha)^2 = + 1.
\]

% ----- printed p. 434 (PDF 460) -----
The foregoing instances of the calculation of $x$, $y$ in the case of the numbers 209 and 173 suggest a table which may be regarded as an extended form of Degen's tables; viz.\ such a table, from $a = 2$ to $a = 99$, is as follows:

\begin{center}
\small\textsc{Specimen of Extended Form of Table in regard to the Pellian Equation.}
\end{center}

\begin{center}
\renewcommand{\arraystretch}{1.05}
\begin{tabular}{|c|r|r|r|}
\hline
$a$ & $y$ & $x$ & $y^2 - ax^2$ \\
\hline
2 & 1 & 0 & $+1$ \\
 & 1 & 1 & $-1$ \\
(2) & 3 & 2 & $+1$ \\
\hline
3 & 1 & 0 & $+1$ \\
 & 1 & 1 & $-2$ \\
(1) & 2 & 1 & $+1$ \\
2 & & & \\
\hline
5 & 2 & 1 & $+1$ \\
 & 2 & 1 & $-1$ \\
(4) & 9 & 4 & $+1$ \\
4 & & & \\
\hline
6 & 2 & 1 & $+1$ \\
 & 2 & 1 & $-2$ \\
(2) & 5 & 2 & $+1$ \\
4 & & & \\
\hline
7 & 2 & 1 & $+1$ \\
1 & 3 & 1 & $-3$ \\
(1) & 5 & 2 & $+2$ \\
1 & 8 & 3 & $-3$ \\
4 & & & $+1$ \\
\hline
8 & 2 & 1 & $+1$ \\
 & 3 & 1 & $-4$ \\
(1) & 17 & 6 & $+1$ \\
4 & & & \\
\hline
10 & 3 & 1 & $+1$ \\
 & 3 & 1 & $-1$ \\
(6) & 19 & 6 & $+1$ \\
6 & & & \\
\hline
11 & 3 & 1 & $+1$ \\
 & 3 & 1 & $-2$ \\
(3) & 10 & 3 & $+1$ \\
6 & & & \\
\hline
12 & 3 & 1 & $+1$ \\
 & 3 & 1 & $-3$ \\
(2) & 7 & 2 & $+1$ \\
6 & & & \\
\hline
13 & 3 & 1 & $+1$ \\
 & 4 & 1 & $-3$ \\
1 & 7 & 2 & $-5$ \\
1 & 11 & 3 & $+4$ \\
1 & 18 & 5 & $+1$ \\
(6) & & & \\
\hline
14 & 3 & 1 & $+1$ \\
 & 4 & 1 & $-1$ \\
(2) & 11 & 3 & $-2$ \\
6 & 15 & 4 & $+1$ \\
\hline
15 & 3 & 1 & $+1$ \\
(1) & 4 & 1 & $+1$ \\
6 & & & \\
\hline
17 & 4 & 1 & $+1$ \\
(8) & 8 & 2 & $-1$ \\
8 & 33 & 8 & $+1$ \\
\hline
18 & 4 & 1 & $+1$ \\
(4) & 9 & 2 & $-2$ \\
8 & 17 & 4 & $+1$ \\
\hline
19 & 4 & 1 & $+1$ \\
 & 4 & 1 & $-3$ \\
(3) & 9 & 2 & $-3$ \\
1 & 13 & 3 & $+5$ \\
2 & 48 & 11 & $-1$ \\
8 & 61 & 14 & $+4$ \\
1 & 170 & 39 & $+1$ \\
\hline
20 & 4 & 1 & $+1$ \\
(2) & 9 & 2 & $-1$ \\
8 & 49 & 11 & $+1$ \\
\hline
\end{tabular}
\end{center}

% ----- printed p. 435 (PDF 461) -----
\begin{center}
\small\textsc{Specimen of Extended Form of Pellian Equation Table---continued.}
\end{center}

\begin{center}
\renewcommand{\arraystretch}{1.05}
\begin{tabular}{|c|r|r|r|}
\hline
$a$ & $y$ & $x$ & $y^2 - ax^2$ \\
\hline
21 & 4 & 1 & $+1$ \\
 & 4 & 1 & $-3$ \\
1 & 5 & 1 & $+4$ \\
1 & 9 & 2 & $-5$ \\
(2) & 23 & 5 & $+4$ \\
1 & 32 & 7 & $-5$ \\
1 & 55 & 12 & $+1$ \\
8 & & & \\
\hline
22 & 4 & 1 & $+1$ \\
1 & 4 & 1 & $-2$ \\
2 & 5 & 1 & $+3$ \\
(4) & 14 & 3 & $+3$ \\
2 & 61 & 13 & $+3$ \\
1 & 136 & 29 & $-6$ \\
8 & 197 & 42 & $+1$ \\
\hline
23 & 4 & 1 & $+1$ \\
1 & 4 & 1 & $-7$ \\
(3) & 5 & 1 & $+2$ \\
1 & 19 & 4 & $-7$ \\
8 & 24 & 5 & $+1$ \\
\hline
24 & 4 & 1 & $+1$ \\
(1) & 4 & 1 & $-1$ \\
8 & 5 & 1 & $+1$ \\
\hline
26 & 5 & 1 & $+1$ \\
(10) & 5 & 1 & $-1$ \\
10 & 51 & 10 & $+1$ \\
\hline
27 & 5 & 1 & $+1$ \\
(5) & 5 & 1 & $-2$ \\
10 & 26 & 5 & $+1$ \\
\hline
28 & 5 & 1 & $+1$ \\
3 & 5 & 1 & $-3$ \\
(2) & 16 & 3 & $+4$ \\
3 & 37 & 7 & $-3$ \\
10 & 127 & 24 & $+1$ \\
\hline
29 & 5 & 1 & $+1$ \\
2 & 5 & 1 & $-4$ \\
(1) & 11 & 2 & $+5$ \\
2 & 16 & 3 & $-5$ \\
10 & 27 & 5 & $+4$ \\
 & 70 & 13 & $-1$ \\
 & 727 & 135 & $+1$ \\
\hline
30 & 5 & 1 & $+1$ \\
(2) & 5 & 1 & $-5$ \\
10 & 11 & 2 & $+1$ \\
\hline
31 & 5 & 1 & $+1$ \\
1 & 5 & 1 & $-6$ \\
1 & 6 & 1 & $+5$ \\
3 & 11 & 2 & $-3$ \\
(5) & 39 & 7 & $+2$ \\
3 & 206 & 37 & $-3$ \\
1 & 657 & 118 & $+5$ \\
1 & 863 & 155 & $-6$ \\
10 & 1520 & 273 & $+1$ \\
\hline
32 & 5 & 1 & $+1$ \\
1 & 5 & 1 & $-7$ \\
(1) & 6 & 1 & $+4$ \\
10 & 17 & 3 & $+1$ \\
\hline
33 & 5 & 1 & $+1$ \\
(2) & 6 & 1 & $-3$ \\
10 & 23 & 4 & $+1$ \\
\hline
34 & 5 & 1 & $+1$ \\
(4) & 6 & 1 & $-2$ \\
10 & 35 & 6 & $+1$ \\
\hline
\end{tabular}
\end{center}

% ----- printed p. 436 (PDF 462) -----
\begin{center}
\small\textsc{Specimen of Extended Form of Pellian Equation Table---continued.}
\end{center}

\begin{center}
\renewcommand{\arraystretch}{1.05}
\begin{tabular}{|c|r|r|r|}
\hline
$a$ & $y$ & $x$ & $y^2 - ax^2$ \\
\hline
35 & 5 & 1 & $+1$ \\
(1) & 6 & 1 & $-1$ \\
10 & & & $+1$ \\
\hline
37 & 6 & 1 & $+1$ \\
(12) & & & \\
12 & 73 & 12 & $+1$ \\
\hline
38 & 6 & 1 & $+1$ \\
(6) & & & \\
12 & 37 & 6 & $+1$ \\
\hline
39 & 6 & 1 & $+1$ \\
(4) & & & \\
12 & 25 & 4 & $+1$ \\
\hline
40 & 6 & 1 & $+1$ \\
(3) & 6 & 1 & $-4$ \\
12 & 19 & 3 & $+1$ \\
\hline
41 & 6 & 1 & $+1$ \\
(2) & 6 & 1 & $-5$ \\
12 & 13 & 2 & $+5$ \\
 & 32 & 5 & $+1$ \\
\hline
42 & 6 & 1 & $+1$ \\
(2) & 6 & 1 & $-6$ \\
12 & 13 & 2 & $+1$ \\
\hline
43 & 6 & 1 & $+1$ \\
1 & 6 & 1 & $-7$ \\
1 & 7 & 1 & $+6$ \\
3 & 13 & 2 & $-3$ \\
1 & 46 & 7 & $+9$ \\
(5) & 59 & 9 & $-2$ \\
1 & 341 & 52 & $+9$ \\
3 & 400 & 61 & $-7$ \\
1 & 1541 & 235 & $+6$ \\
1 & 1941 & 296 & $-7$ \\
12 & 3482 & 531 & $+1$ \\
\hline
44 & 6 & 1 & $+1$ \\
 & 6 & 1 & $-8$ \\
1 & 7 & 1 & $+5$ \\
1 & 13 & 2 & $-7$ \\
1 & 20 & 3 & $+4$ \\
(2) & 53 & 8 & $-5$ \\
1 & 73 & 11 & $+8$ \\
1 & 126 & 19 & $-5$ \\
12 & 199 & 30 & $+1$ \\
\hline
45 & 6 & 1 & $+1$ \\
 & 6 & 1 & $-9$ \\
(2) & 7 & 1 & $+4$ \\
2 & 20 & 3 & $-9$ \\
12 & 47 & 7 & $+4$ \\
 & 114 & 17 & $-9$ \\
 & 161 & 24 & $+4$ \\
\hline
46 & 6 & 1 & $+1$ \\
1 & 6 & 1 & $-10$ \\
3 & 7 & 1 & $+3$ \\
1 & 27 & 4 & $-7$ \\
1 & 34 & 5 & $+6$ \\
2 & 61 & 9 & $-5$ \\
6 & 156 & 23 & $+2$ \\
1 & 997 & 147 & $-7$ \\
1 & 2150 & 317 & $-7$ \\
3 & 3147 & 464 & $+6$ \\
1 & 5297 & 781 & $-11$ \\
12 & 24335 & 3588 & $+1$ \\
\hline
47 & 6 & 1 & $+1$ \\
1 & 6 & 1 & $-11$ \\
(5) & 417 & 60 & $-12$ \\
12 & 48 & 7 & $+2$ \\
\hline
48 & 6 & 1 & $+1$ \\
(1) & 6 & 1 & $+1$ \\
12 & 67 & 1 & $+1$ \\
\hline
\end{tabular}
\end{center}

% ----- printed p. 437 (PDF 463) -----
\begin{center}
\small\textsc{Specimen of Extended Form of Pellian Equation Table---continued.}
\end{center}

\begin{center}
\renewcommand{\arraystretch}{1.05}
\begin{tabular}{|c|r|r|r|}
\hline
$a$ & $y$ & $x$ & $y^2 - ax^2$ \\
\hline
50 & 7 & 1 & $+1$ \\
(14) & 7 & 1 & $-1$ \\
14 & 99 & 14 & $+1$ \\
\hline
51 & 7 & 1 & $+1$ \\
(7) & 7 & 1 & $-2$ \\
14 & 50 & 7 & $+1$ \\
\hline
52 & 7 & 1 & $+1$ \\
4 & 7 & 1 & $-3$ \\
1 & 29 & 4 & $-4$ \\
(2) & 36 & 5 & $+9$ \\
1 & 101 & 14 & $-3$ \\
4 & 137 & 19 & $+9$ \\
14 & 649 & 90 & $+1$ \\
\hline
53 & 7 & 1 & $+1$ \\
3 & 7 & 1 & $-4$ \\
(1) & 22 & 3 & $+7$ \\
3 & 29 & 4 & $-7$ \\
14 & 51 & 7 & $+4$ \\
 & 182 & 25 & $-5$ \\
\hline
54 & 7 & 1 & $+1$ \\
2 & 7 & 1 & $-5$ \\
1 & 15 & 2 & $+2$ \\
(6) & 22 & 3 & $-3$ \\
1 & 147 & 20 & $+9$ \\
2 & 169 & 23 & $-1$ \\
14 & 485 & 66 & $+9$ \\
\hline
55 & 7 & 1 & $+1$ \\
2 & 7 & 1 & $-6$ \\
(2) & 15 & 2 & $+5$ \\
2 & 37 & 5 & $-6$ \\
14 & 89 & 12 & $+1$ \\
\hline
56 & 7 & 1 & $+1$ \\
(2) & 7 & 1 & $-7$ \\
14 & 15 & 2 & $+1$ \\
\hline
57 & 7 & 1 & $+1$ \\
1 & 7 & 1 & $-8$ \\
1 & 8 & 1 & $+7$ \\
(4) & 15 & 2 & $-3$ \\
1 & 68 & 9 & $+7$ \\
1 & 83 & 11 & $-8$ \\
14 & 151 & 20 & $+1$ \\
\hline
58 & 7 & 1 & $+1$ \\
1 & 7 & 1 & $-9$ \\
1 & 8 & 1 & $+6$ \\
(1) & 15 & 2 & $-7$ \\
1 & 23 & 3 & $+9$ \\
1 & 38 & 5 & $-10$ \\
14 & 99 & 13 & $+1$ \\
\hline
59 & 7 & 1 & $+1$ \\
1 & 7 & 1 & $-10$ \\
2 & 8 & 1 & $+5$ \\
(7) & 23 & 3 & $-5$ \\
2 & 169 & 22 & $-11$ \\
1 & 361 & 47 & $-10$ \\
14 & 530 & 69 & $+1$ \\
\hline
60 & 7 & 1 & $+1$ \\
(2) & 7 & 1 & $-11$ \\
14 & 8 & 1 & $+4$ \\
 & 238 & 31 & $+1$ \\
\hline
61 & 7 & 1 & $+1$ \\
1 & 7 & 1 & $-12$ \\
4 & 8 & 1 & $+3$ \\
3 & 39 & 5 & $-4$ \\
1 & 125 & 16 & $+9$ \\
(2) & 164 & 21 & $-5$ \\
 & 453 & 58 & $+5$ \\
1 & 1070 & 137 & $-9$ \\
\hline
\end{tabular}
\end{center}
% ----- printed p. 438 (PDF 464) -----
\begin{center}
\small\textsc{Specimen of Extended Form of Pellian Equation Table---continued.}
\end{center}

\begin{center}
\renewcommand{\arraystretch}{1.05}
\begin{tabular}{|c|r|r|r|}
\hline
$a$ & $y$ & $x$ & $y^2 - ax^2$ \\
\hline
 & 1523 & 195 & $+4$ \\
 & 5639 & 722 & $-3$ \\
1 & 24079 & 3083 & $+12$ \\
14 & 29718 & 3805 & $-13$ \\
\hline
62 & 7 & 1 & $+1$ \\
1 & 7 & 1 & $-13$ \\
(6) & 8 & 1 & $+2$ \\
1 & 55 & 7 & $-7$ \\
14 & 63 & 8 & $+1$ \\
\hline
63 & 7 & 1 & $+1$ \\
(1) & 7 & 1 & $-14$ \\
14 & 8 & 1 & $+1$ \\
\hline
65 & 8 & 1 & $+1$ \\
(16) & 8 & 1 & $-1$ \\
16 & 129 & 16 & $+1$ \\
\hline
66 & 8 & 1 & $+1$ \\
(8) & 8 & 1 & $-2$ \\
16 & 65 & 8 & $+1$ \\
\hline
67 & 8 & 1 & $+1$ \\
5 & 8 & 1 & $-3$ \\
2 & 41 & 5 & $-6$ \\
1 & 90 & 11 & $+7$ \\
1 & 131 & 16 & $-9$ \\
(7) & 221 & 27 & $+9$ \\
1 & 1678 & 205 & $+9$ \\
1 & 1899 & 232 & $+7$ \\
2 & 3577 & 437 & $-6$ \\
5 & 9053 & 1106 & $+3$ \\
16 & 48842 & 5967 & $+1$ \\
\hline
68 & 8 & 1 & $+1$ \\
(4) & 8 & 1 & $-4$ \\
16 & 33 & 4 & $+1$ \\
\hline
69 & 8 & 1 & $+1$ \\
3 & 8 & 1 & $-5$ \\
1 & 25 & 3 & $-5$ \\
1 & 83 & 10 & $+4$ \\
(4) & 108 & 13 & $-5$ \\
1 & 515 & 62 & $+4$ \\
3 & 623 & 75 & $-5$ \\
3 & 2384 & 297 & $-6$ \\
16 & 7775 & 936 & $+1$ \\
\hline
70 & 8 & 1 & $+1$ \\
2 & 8 & 1 & $-6$ \\
1 & 17 & 2 & $+3$ \\
(2) & 25 & 3 & $-7$ \\
1 & 67 & 8 & $+9$ \\
2 & 92 & 11 & $-6$ \\
16 & 251 & 30 & $-11$ \\
\hline
71 & 8 & 1 & $+1$ \\
2 & 8 & 1 & $-7$ \\
2 & 17 & 2 & $+2$ \\
1 & 42 & 5 & $-11$ \\
(7) & 59 & 7 & $+2$ \\
2 & 455 & 54 & $-9$ \\
2 & 514 & 61 & $+5$ \\
16 & 3480 & 413 & $+1$ \\
\hline
72 & 8 & 1 & $+1$ \\
(2) & 8 & 1 & $-8$ \\
16 & 17 & 2 & $+1$ \\
\hline
73 & 8 & 1 & $+1$ \\
1 & 8 & 1 & $-9$ \\
1 & 9 & 1 & $+8$ \\
(5) & 17 & 2 & $-3$ \\
1 & 94 & 11 & $+8$ \\
1 & 111 & 13 & $-9$ \\
16 & 1068 & 125 & $+9$ \\
\hline
74 & 8 & 1 & $+1$ \\
1 & 8 & 1 & $-10$ \\
\hline
\end{tabular}
\end{center}
% ----- printed p. 439 (PDF 465) -----
\begin{center}
\small\textsc{Specimen of Extended Form of Pellian Equation Table---continued.}
\end{center}

\begin{center}
\renewcommand{\arraystretch}{1.05}
\begin{tabular}{|c|r|r|r|}
\hline
$a$ & $y$ & $x$ & $y^2 - ax^2$ \\
\hline
(1) & 9 & 1 & $+7$ \\
1 & 17 & 2 & $-7$ \\
16 & 26 & 3 & $+10$ \\
 & 43 & 5 & $-7$ \\
\hline
75 & 8 & 1 & $+1$ \\
(1) & 8 & 1 & $-11$ \\
1 & 9 & 1 & $+6$ \\
16 & 17 & 2 & $-11$ \\
 & 26 & 3 & $+1$ \\
\hline
76 & 8 & 1 & $+1$ \\
1 & 8 & 1 & $-12$ \\
1 & 9 & 1 & $+5$ \\
2 & 26 & 3 & $-8$ \\
1 & 35 & 4 & $+9$ \\
1 & 61 & 7 & $-3$ \\
5 & 340 & 39 & $+4$ \\
(4) & 1421 & 163 & $-3$ \\
5 & 7445 & 854 & $-12$ \\
1 & 8866 & 1017 & $+5$ \\
1 & 16311 & 1871 & $+1$ \\
2 & 41488 & 4759 & $+1$ \\
16 & 57799 & 6630 & $+1$ \\
\hline
77 & 8 & 1 & $+1$ \\
3 & 9 & 1 & $+4$ \\
(2) & 35 & 4 & $-13$ \\
3 & 79 & 9 & $+4$ \\
1 & 272 & 31 & $-1$ \\
16 & 351 & 40 & $+1$ \\
\hline
78 & 8 & 1 & $+1$ \\
1 & 8 & 1 & $-14$ \\
(4) & 9 & 1 & $+3$ \\
1 & 44 & 5 & $-14$ \\
16 & 53 & 6 & $+1$ \\
\hline
79 & 8 & 1 & $+1$ \\
1 & 8 & 1 & $-15$ \\
(7) & 9 & 1 & $+2$ \\
1 & 71 & 8 & $-15$ \\
16 & 80 & 9 & $+1$ \\
\hline
80 & 8 & 1 & $+1$ \\
(1) & 8 & 1 & $-16$ \\
16 & 9 & 1 & $+1$ \\
\hline
82 & 9 & 1 & $+1$ \\
(18) & 9 & 1 & $-1$ \\
18 & 163 & 18 & $+1$ \\
\hline
83 & 9 & 1 & $+1$ \\
(9) & 9 & 1 & $-2$ \\
18 & 82 & 9 & $+1$ \\
\hline
84 & 9 & 1 & $+1$ \\
(6) & 9 & 1 & $-3$ \\
18 & 55 & 6 & $+1$ \\
\hline
85 & 9 & 1 & $+1$ \\
4 & 9 & 1 & $-4$ \\
(1) & 37 & 4 & $+9$ \\
1 & 46 & 5 & $-5$ \\
4 & 83 & 9 & $-7$ \\
18 & 378 & 41 & $+4$ \\
\hline
86 & 9 & 1 & $+1$ \\
3 & 9 & 1 & $-5$ \\
1 & 28 & 3 & $-12$ \\
1 & 37 & 4 & $+5$ \\
1 & 65 & 7 & $-7$ \\
(8) & 102 & 11 & $-5$ \\
1 & 881 & 95 & $+11$ \\
1 & 983 & 106 & $+11$ \\
1 & 1864 & 201 & $-5$ \\
3 & 2847 & 307 & $+10$ \\
18 & 10405 & 1122 & $+1$ \\
\hline
\end{tabular}
\end{center}

% ----- printed p. 440 (PDF 466) -----
\begin{center}
\small\textsc{Specimen of Extended Form of Pellian Equation Table---continued.}
\end{center}

\begin{center}
\renewcommand{\arraystretch}{1.05}
\begin{tabular}{|c|r|r|r|}
\hline
$a$ & $y$ & $x$ & $y^2 - ax^2$ \\
\hline
87 & 9 & 1 & $+1$ \\
(3) & 9 & 1 & $-6$ \\
18 & 28 & 3 & $+1$ \\
\hline
88 & 9 & 1 & $+1$ \\
2 & 9 & 1 & $-7$ \\
1 & 19 & 2 & $+9$ \\
(1) & 28 & 3 & $-8$ \\
1 & 47 & 5 & $+9$ \\
2 & 75 & 8 & $-7$ \\
18 & 197 & 21 & $+1$ \\
\hline
89 & 9 & 1 & $+1$ \\
2 & 9 & 1 & $-8$ \\
(3) & 19 & 2 & $+5$ \\
2 & 66 & 7 & $-5$ \\
18 & 217 & 23 & $+8$ \\
 & 500 & 53 & $+1$ \\
\hline
90 & 9 & 1 & $+1$ \\
(2) & 9 & 1 & $-9$ \\
18 & 19 & 2 & $+1$ \\
\hline
91 & 9 & 1 & $+1$ \\
1 & 9 & 1 & $-10$ \\
1 & 10 & 1 & $+9$ \\
5 & 19 & 2 & $-3$ \\
(1) & 105 & 11 & $-6$ \\
5 & 124 & 13 & $+5$ \\
1 & 725 & 76 & $-9$ \\
1 & 849 & 89 & $+10$ \\
18 & 1574 & 165 & $+1$ \\
\hline
92 & 9 & 1 & $+1$ \\
1 & 9 & 1 & $-11$ \\
1 & 10 & 1 & $+8$ \\
2 & 19 & 2 & $-7$ \\
(4) & 48 & 5 & $-11$ \\
2 & 211 & 22 & $+5$ \\
1 & 470 & 49 & $-7$ \\
1 & 681 & 71 & $+8$ \\
18 & 1151 & 120 & $-11$ \\
\hline
93 & 9 & 1 & $+1$ \\
1 & 9 & 1 & $-12$ \\
1 & 10 & 1 & $+7$ \\
1 & 19 & 2 & $-11$ \\
4 & 29 & 3 & $+4$ \\
(6) & 135 & 14 & $-9$ \\
4 & 839 & 87 & $-4$ \\
1 & 3491 & 362 & $+5$ \\
1 & 4330 & 449 & $-11$ \\
1 & 7821 & 811 & $+5$ \\
18 & 12151 & 1260 & $-12$ \\
\hline
94 & 9 & 1 & $+1$ \\
1 & 9 & 1 & $-13$ \\
2 & 10 & 1 & $+6$ \\
3 & 29 & 3 & $-13$ \\
1 & 97 & 10 & $+9$ \\
1 & 126 & 13 & $-15$ \\
5 & 223 & 23 & $+2$ \\
1 & 1241 & 128 & $-15$ \\
(8) & 1464 & 151 & $+2$ \\
1 & 12953 & 1336 & $+5$ \\
5 & 14417 & 1487 & $-3$ \\
1 & 85038 & 8771 & $+5$ \\
1 & 99455 & 10258 & $-14$ \\
3 & 184493 & 19029 & $+6$ \\
2 & 652934 & 67345 & $-14$ \\
1 & 1490361 & 153719 & $+1$ \\
18 & 2143295 & 221064 & $+1$ \\
\hline
95 & 9 & 1 & $+1$ \\
(2) & 9 & 1 & $-14$ \\
18 & 29 & 3 & $+5$ \\
 & 39 & 4 & $-15$ \\
\hline
96 & 9 & 1 & $+1$ \\
(3) & 9 & 1 & $-15$ \\
18 & 10 & 1 & $+4$ \\
 & 39 & 4 & $+1$ \\
\hline
\end{tabular}
\end{center}
% ----- printed p. 441 (PDF 467) -----
\begin{center}
\small\textsc{Specimen of Extended Form of Pellian Equation Table---continued.}
\end{center}

\begin{center}
\renewcommand{\arraystretch}{1.05}
\begin{tabular}{|c|r|r|r|}
\hline
$a$ & $y$ & $x$ & $y^2 - ax^2$ \\
\hline
97 & 9 & 1 & $+1$ \\
1 & 9 & 1 & $-16$ \\
5 & 10 & 1 & $+3$ \\
1 & 59 & 6 & $-13$ \\
1 & 69 & 7 & $+8$ \\
(1) & 128 & 13 & $-9$ \\
1 & 197 & 20 & $+9$ \\
1 & 325 & 33 & $-8$ \\
5 & 522 & 53 & $+13$ \\
1 & 847 & 86 & $-3$ \\
1 & 4757 & 483 & $+16$ \\
18 & 5604 & 569 & $+1$ \\
\hline
98 & 9 & 1 & $+1$ \\
(8) & 9 & 1 & $-17$ \\
18 & 10 & 1 & $+2$ \\
 & 89 & 9 & $-17$ \\
 & 99 & 10 & $+1$ \\
\hline
99 & 9 & 1 & $+1$ \\
(1) & 9 & 1 & $-18$ \\
18 & 10 & 1 & $+1$ \\
\hline
\end{tabular}
\end{center}

The meaning hardly requires explanation; for each number $a$, we have a series of pairs of increasing numbers, $y$, $x$, satisfying a series of equations $y^2 = ax^2 \pm b$; thus
\[
a = 14
\]
\begin{center}
\begin{tabular}{ccr}
$y$ & $x$ & $y^2 - ax^2$ \\
\hline
1 & 0 & $1 - 14 \cdot 0 = 1$, \\
3 & 1 & $9 - 14 \cdot 1 = -5$, \\
4 & 1 & $16 - 14 \cdot 1 = +2$, \\
11 & 3 & $121 - 14 \cdot 9 = -5$, \\
15 & 4 & $225 - 14 \cdot 16 = +1$. \\
\end{tabular}
\end{center}

The following table, calculated under the superintendence of the Committee, extends from $a = 1001$ to $a = 1500$ (square numbers omitted); it is (with slight typographical variations) nearly but not exactly in the form of Degen's Table I., the chief difference being that for a number $a$ having a double middle term, or of the form $\alpha^2 + 1$ (such number being further distinguished by an asterisk), the $x$, $y$ entered in the table are the solutions, not of the equation $y^2 = ax^2 + 1$, but of the equation $y^2 = ax^2 - 1$. As remarked above, if we have $y^2 = ax^2 - 1$, then writing $y_1 = 2y^2 + 1$ and $x_1 = 2xy$, we obtain $y_1^2 = ax_1^2 + 1$.

Moreover, for each value of $a$, in the first line, the first term, which is the integer part of $\sqrt{a}$, is separated from the others by a semicolon, and the 1, which is the corresponding first term of the second line, is omitted.

% ----- printed p. 442 (PDF 468) -----
The calculations were made by C.~E. Bickmore, M.A., of New College, Oxford: his values for $x$ and $y$ have been revised as presently mentioned, but it has been assumed that his values for the periods and subsidiary numbers (forming the first and second lines of each division of the table) are accurate; in fact, any error therein would cause the resulting values of $x$ and $y$ to be wildly erroneous; but (except in a single instance which was accounted for) the errors in $x$ and $y$ were in every case in a single figure or two or three figures only.

The values of $x$ and $y$ were in every case examined by substitution in the equation ($y^2 = ax^2 + 1$, or $y^2 = ax^2 - 1$, as the case may be), which should be satisfied by them. These verifications were for the most part made by A. Graham, M.A., of the Observatory, Cambridge. As already mentioned, some errors were detected, and these have been, of course, corrected. The values of $x$, $y$ given in the table thus satisfy in every case the proper equation $y^2 = ax^2 + 1$, or $y^2 = ax^2 - 1$; on the ground above referred to, it is believed that the periods and subsidiary numbers are also accurate.

It may be remarked, in regard to the verification of the equation $y^2 = ax^2 \pm 1$ for large values of $a$ and $y$, it is in practice easier and safer to calculate $ax^2 \pm 1$, and then to compare the square root thereof with the given value of $y$, than to further calculate the value of $y^2$.

% ----- printed p. 443 (PDF 469) -----
\begin{center}
\small\textsc{The Table 1001 to 1500.}
\end{center}

\noindent The following table lists, for each integer $a$ from 1001 to 1500 (perfect squares omitted), the continued-fraction expansion of $\sqrt{a}$, the auxiliary period-numbers, and the resulting fundamental solution $(x,y)$ of the Pellian equation. Entries marked by an asterisk $\asterism$ correspond to a double middle term; in those cases the tabulated $(x,y)$ satisfy $y^2 = ax^2 - 1$. A semicolon separates the integer part from the period; parenthesised numbers denote middle terms.

\medskip
\renewcommand{\arraystretch}{1.15}
\begin{longtable}{|r|p{0.55\textwidth}|r|}
\hline
$a$ & continued-fraction expansion of $\sqrt{a}$; period and subsidiary numbers & $x$ \\
 & & $y$ \\
\hline
\endhead
1001 & $31\,;\ 1, 1, 1, 3, 3, 2, (4)$; \ \ \ \ $40,\ 23,\ 32,\ 16,\ 17,\ 25,\ (13)$ & 33532 \\
 & & 10\,60905 \\
\hline
1002 & $31\,;\ 1, 1, 1, 8, 2, 1, 1, 1, 3, (10)$; \ \ \ \ $41,\ 22,\ 39,\ 7,\ 23,\ 31,\ 26,\ 33,\ 17,\ (6)$ & 65\,35248 \\
 & & 2068\,69247 \\
\hline
1003 & $31\,;\ 1, 2, (31)$; \ \ \ \ $42,\ 21,\ (2)$ & 285 \\
 & & 9026 \\
\hline
1004 & $31\,;\ 1, 2, 5, 2, 2, 1, 7, 4, 1, 2, 1, 11, 1, (14)$; \ \ \ \ $43,\ 20,\ 11,\ 25,\ 19,\ 41,\ 8,\ 13,\ 40,\ 17,\ 44,\ 5,\ 55,\ (4)$ & 85\,24164\,59730 \\
 & & 2700\,96330\,24199 \\
\hline
1005 & $31\,;\ 1, 2, 2, 1, 5, 15, 1, 2, (12)$; \ \ \ \ $44,\ 19,\ 20,\ 39,\ 11,\ 4,\ 41,\ 21,\ (5)$ & 930\,59568 \\
 & & 29501\,49761 \\
\hline
1006 & $31\,;\ 1, 2, 1, 1, 5, 1, 3, 2, 1, 1, 1, 1, 1, 9, 1, 20, 4, 5, 1, 1, 12, 6, 1, (30)$; \ \ \ \ $45,\ 18,\ 29,\ 33,\ 10,\ 43,\ 15,\ 22,\ 31,\ 27,\ 30,\ 25,\ 37,\ 6,\ 55,\ 3,\ 15,\ 11,\ 30,\ 33,\ 5,\ 9,\ 53,\ (2)$ & 4\,45346\,14025\,55749\,21748 \\
 & & 141\,25267\,56378\,02146\,05455 \\
\hline
1007 & $31\,;\ 1, 2, (1)$; \ \ \ \ $46,\ 17,\ (38)$ & 127 \\
 & & 17 \\
\hline
1008 & $31\,;\ 1, (2)$; \ \ \ \ $47,\ (16)$ & 4 \\
 & & 127 \\
\hline
1009$\asterism$ & $31\,;\ 1, (3, 3)$; \ \ \ \ $48,\ (15, 15)$ & 17 \\
 & & 540 \\
\hline
1010$\asterism$ & $31\,;\ 1, 3, (1, 1)$; \ \ \ \ $49,\ 14,\ (31, 31)$ & 265 \\
 & & 1303 \\
\hline
1011 & $31\,;\ 1, 3, 1, (9)$; \ \ \ \ $50,\ 13,\ 47,\ (6)$ & 41 \\
 & & 8426 \\
\hline
1012 & $31\,;\ 1, 4, 3, 6, 1, 3, 8, 1, (4)$; \ \ \ \ $51,\ 12,\ 19,\ 9,\ 43,\ 16,\ 7,\ 48,\ (11)$ & 1013\,02110 \\
 & & 32226\,17399 \\
\hline
1013$\asterism$ & $31\,;\ 1, 4, 1, 4, 15, 1, 2, (2, 2)$; \ \ \ \ $52,\ 11,\ 44,\ 13,\ 4,\ 43,\ 19,\ (23, 23)$ & 123\,52985 \\
 & & 3931\,66618 \\
\hline
1014 & $31\,;\ 1, 5, 2, 1, 1, 1, 1, (20)$; \ \ \ \ $53,\ 10,\ 23,\ 30,\ 29,\ 25,\ 38,\ (3)$ & 1\,46266 \\
 & & 46\,56965 \\
\hline
1015 & $31\,;\ 1, 6, 10, (2)$; \ \ \ \ $54,\ 9,\ 6,\ (29)$ & 11076 \\
 & & 3\,52871 \\
\hline
1016 & $31\,;\ 1, (6)$; \ \ \ \ $55,\ (8)$ & 2558 \\
 & & 9\,09655\,84992 \\
\hline
1017 & $31\,;\ 1, 8, 1, 6, 1, 1, 4, 2, 1, 5, 5, (1)$; \ \ \ \ $56,\ 7,\ 8,\ 49,\ 9,\ 13,\ 41,\ 16,\ 11,\ 31,\ 32,\ (9)$ & 290\,09322\,97217 \\
 & & 27\,28333 \\
\hline
1018$\asterism$ & $31\,;\ 1, 9, 1, 1, 1, 6, 2, (3, 3)$; \ \ \ \ $57,\ 6,\ 39,\ 23,\ 38,\ 9,\ 26,\ (17, 17)$ & 870\,50499 \\
 & & 19\,07764\,36539 \\
\hline
1019 & $31\,;\ 1, 11, 1, 3, 1, 1, 1, 3, 8, 1, 5, 2, (31)$; \ \ \ \ $58,\ 5,\ 47,\ 14,\ 35,\ 25,\ 34,\ 17,\ 7,\ 49,\ 10,\ 29,\ (2)$ & 608\,99233\,21730 \\
 & & 16 \\
\hline
1020 & $31\,;\ 1, (14)$; \ \ \ \ $59,\ (4)$ & 511 \\
 & & 98\,65001\,29666\,69564\,06909 \\
\hline
1021$\asterism$ & $31\,;\ 1, 20, 3, 6, 1, 3, 2, 1, 1, 12, 5, 4, 15, 1, 2, 1, 4, 1, 1, 2, 1, 1, 1, (5, 5)$; \ \ \ \ $60,\ 3,\ 20,\ 9,\ 44,\ 15,\ 23,\ 27,\ 36,\ 5,\ 12,\ 15,\ 4,\ 45,\ 17,\ 41,\ 12,\ 33,\ 29,\ 20,\ 33,\ 25,\ 36,\ (11, 11)$ & 3152\,17280\,37258\,48825\,15030 \\
 & & 1023 \\
\hline
\end{longtable}

% ----- printed p. 444 (PDF 470) -----
\begin{center}
\small\textsc{Table 1001 to 1500---continued.}
\end{center}

\renewcommand{\arraystretch}{1.15}
\begin{longtable}{|r|p{0.55\textwidth}|r|}
\hline
$a$ & continued-fraction expansion of $\sqrt{a}$; period and subsidiary numbers & $x$ \\
 & & $y$ \\
\hline
\endhead
1022 & $31\,;\ 1, (30)$; \ \ \ \ $61,\ (2)$ & 32 \\
 & & 1023 \\
\hline
1023 & $31\,;\ (1)$; \ \ \ \ $(62)$ & 1 \\
 & & 32 \\
\hline
1025$\asterism$ & $32\,;\ (64)$; \ \ \ \ $(1)$ & 1 \\
 & & 32 \\
\hline
1026 & $32\,;\ (32)$; \ \ \ \ $(2)$ & 32 \\
 & & 1025 \\
\hline
1027 & $32\,;\ 21, 1, 1, 6, 2, (4)$; \ \ \ \ $3,\ 22,\ 39,\ 9,\ 27,\ (13)$ & 41\,54868 \\
 & & 1331\,50393 \\
\hline
1028 & $32\,;\ (16)$; \ \ \ \ $(4)$ & 16 \\
 & & 513 \\
\hline
1029 & $32\,;\ 12, 1, 4, 2, 2, 1, 3, 15, 1, 3, 2, 1, (20)$; \ \ \ \ $5,\ 49,\ 12,\ 25,\ 20,\ 37,\ 17,\ 4,\ 47,\ 15,\ 20,\ 43,\ (3)$ & 303\,57068\,95884 \\
 & & 9737\,94964\,66615 \\
\hline
1030 & $32\,;\ 10, 1, 2, 6, 1, 3, 1, 2, 1, 1, 2, 2, (12)$; \ \ \ \ $6,\ 41,\ 21,\ 9,\ 45,\ 14,\ 39,\ 19,\ 31,\ 30,\ 21,\ 26,\ (5)$ & 2\,75531\,87434 \\
 & & 88\,43070\,11291 \\
\hline
1031 & $32\,;\ 9, 6, 3, 4, 1, 1, 1, 12, 4, 1, 1, (31)$; \ \ \ \ $7,\ 10,\ 19,\ 13,\ 35,\ 26,\ 31,\ 25,\ 38,\ 5,\ 14,\ 29,\ 35,\ (2)$ & 2029\,75370\,82877 \\
 & & 65173\,74486\,64200 \\
\hline
1032 & $32\,;\ (8)$; \ \ \ \ $(8)$ & 2578 \\
 & & 82780 \\
\hline
1033$\asterism$ & $32\,;\ 7, 7, 1, 8, 3, 3, 1, 2, 3, 2, 2, 1, 1, 1, 2, (21, 21)$; \ \ \ \ $9,\ 8,\ 51,\ 7,\ 19,\ 16,\ 37,\ 21,\ 17,\ 24,\ 21,\ 32,\ 27,\ 31,\ 24,\ (3, 3)$ & 5\,81389\,30460\,24093 \\
 & & 186\,86036\,75961\,74196 \\
\hline
1034 & $32\,;\ 6, 2, (2)$; \ \ \ \ $10,\ 25,\ (22)$ & 494 \\
 & & 15885 \\
\hline
1035 & $32\,;\ 11, 5, (46)$; \ \ \ \ $(1)$ & 35 \\
 & & 1126 \\
\hline
1036 & $32\,;\ 12, 5, 2, 1, 6, (2)$; \ \ \ \ $21,\ 40,\ 9,\ (28)$ & 26322 \\
 & & 8\,47225 \\
\hline
1037$\asterism$ & $32\,;\ 4, 1, 15, (3, 3)$; \ \ \ \ $13,\ 49,\ 4,\ (19, 19)$ & 64805 \\
 & & 20\,86882 \\
\hline
1038 & $32\,;\ 4, 1, 1, 2, 2, 1, 2, (10)$; \ \ \ \ $14,\ 33,\ 29,\ 22,\ 21,\ 34,\ 23,\ (6)$ & 4\,78126 \\
 & & 154\,04267 \\
\hline
1039 & $32\,;\ 4, 3, 1, 1, 5, 3, 2, 2, 21, 12, 1, 5, 1, 1, 10, 4, 1, 6, 2, 1, 3, 1, 1, 1, 1, 1, 1, (31)$; \ \ \ \ $15,\ 17,\ 30,\ 33,\ 11,\ 18,\ 23,\ 26,\ 3,\ 5,\ 51,\ 10,\ 31,\ 33,\ 6,\ 13,\ 46,\ 9,\ 22,\ 37,\ 15,\ 34,\ 27,\ 29,\ 30,\ 25,\ 39,\ (2)$ & 1\,53006\,74275\,15667\,18643\,06921 \\
 & & 49\,31946\,34979\,07633\,93486\,37520 \\
\hline
1040 & $32\,;\ (4)$; \ \ \ \ $(16)$ & 1294 \\
 & & 32 \\
\hline
1041 & $32\,;\ 3, 1, 3, 1, 1, 4, 2, 2, 7, 1, 1, 1, 12, 3, 1, (20)$; \ \ \ \ $17,\ 40,\ 15,\ 31,\ 32,\ 13,\ 24,\ 25,\ 8,\ 39,\ 23,\ 40,\ 5,\ 16,\ 47,\ (3)$ & 2389\,36879\,43492 \\
 & & 77091\,86499\,27575 \\
\hline
1042$\asterism$ & $32\,;\ 3, (1, 1)$; \ \ \ \ $18,\ (31, 31)$ & 807 \\
 & & 25 \\
\hline
\end{longtable}

% ----- printed p. 445 (PDF 471) -----
\begin{center}
\small\textsc{Table 1001 to 1500---continued.}
\end{center}

\renewcommand{\arraystretch}{1.15}
\begin{longtable}{|r|p{0.55\textwidth}|r|}
\hline
$a$ & continued-fraction expansion of $\sqrt{a}$; period and subsidiary numbers & $x$ \\
 & & $y$ \\
\hline
\endhead
1043 & $32\,;\ 3, 2, 1, 1, 1, 4, (8)$; \ \ \ \ $19,\ 22,\ 31,\ 29,\ 26,\ 37,\ (7)$ & 59\,76927 \\
 & & 5\,17864 \\
\hline
1044 & $32\,;\ 20, 1, 1, 1, 3, 2, 1, 1, (6)$; \ \ \ \ $21,\ 15,\ 36,\ 25,\ 35,\ 16,\ 23,\ 28,\ 35,\ (9)$ & 45940 \\
 & & 1457\,20107 \\
\hline
1045 & $32\,;\ 3, 1, (4)$; \ \ \ \ $22,\ 24,\ 51,\ (11)$ & 14112 \\
 & & 4\,56191 \\
\hline
1046 & $32\,;\ 23, 13, 1, 5, 17, 38, 19, 34, 25, 37, 10, 47, 11, (2)$; \ \ \ \ $5,\ 17,\ 38,\ 19,\ 34,\ 25,\ 37,\ 10,\ 47,\ 11,\ (2)$ & 12\,32778\,42162 \\
 & & 398\,70425 \\
\hline
1047 & $32\,;\ 24, 3, 49, 1, (20)$; \ \ \ \ $14,\ 59,\ (3)$ & 36807 \\
 & & 11\,90060 \\
\hline
1048 & $32\,;\ 2, 4, 37, 1, 2, 6, 1, 4, 1, 1, (7)$; \ \ \ \ $25,\ 32,\ 9,\ 47,\ 12,\ 31,\ 33,\ (8)$ & 1\,42859 \\
 & & 46\,18316 \\
\hline
1049$\asterism$ & $32\,;\ 25, 33, 2, 23, 9, 47, 12, 31, 33, (8)$; \ \ \ \ $32,\ 19,\ 40,\ 13,\ 16,\ 43,\ 11,\ 8,\ 53,\ (5, 5)$ & 1534\,09493 \\
 & & 49654\,25620 \\
\hline
1050 & $32\,;\ 2, 1, 25, (10)$; \ \ \ \ $26,\ 25,\ (6)$ & 270 \\
 & & 8749 \\
\hline
1051 & $32\,;\ 2, 17, 1, 1, 14, 15, 6, 35, 29, 18, 5, 42, 21, 10, 9, 3, 30, 7, 46, 15, 25, 19, 4, 5, (2)$; \ \ \ \ $27,\ 2,\ 1,\ 30,\ 23,\ 4,\ 4,\ 10,\ 1,\ 1,\ 3,\ 12,\ 1,\ 2,\ 6,\ 7,\ 21,\ 2,\ 8,\ 1,\ 3,\ 2,\ 2,\ 1,\ (31)$ & 11935\,20 \\
 & & 2\,94697\,22410\,66655\,86570\,73250 \\
\hline
1052 & $32\,;\ 28, 19, 3, 7, 1, 4, 9, (16)$; \ \ \ \ $3,\ 19,\ 8,\ 35,\ 13,\ 7,\ (4)$ & 3\,00767 \\
 & & 87132\,79112 \\
\hline
1053 & $32\,;\ 29, 2, 3, (4)$; \ \ \ \ $(13)$ & 649 \\
 & & 30 \\
\hline
1054 & $32\,;\ 30, 9, 62, 1, 2, 1, 1, 4, 2, 2, 1, 1, 1, 3, 1, 2, 3, 4, (32)$; \ \ \ \ $2,\ 19,\ 50,\ 33,\ 13,\ 25,\ 21,\ 33,\ 26,\ 35,\ 15,\ 38,\ 21,\ 18,\ 15,\ (2)$ & 10335\,86450\,05416 \\
 & & 3\,35557\,62560\,56895 \\
\hline
1055 & $32\,;\ (12)$; \ \ \ \ $(5)$ & 1689 \\
 & & 52 \\
\hline
1056 & $32\,;\ 1, (5)$; \ \ \ \ $(32)$ & 2 \\
 & & 65 \\
\hline
1057 & $32\,;\ 21, 5, 1, 6, 2, 1, 1, 3, 2, 7, 1, 2, 4, 1, 1, 1, (8)$; \ \ \ \ $33,\ 1,\ 3, 11,\ 48,\ 9,\ 24,\ 29,\ 33,\ 16,\ 27,\ 8,\ 41,\ 21,\ 13,\ 37,\ 24,\ 39,\ (7)$ & 32992 \\
 & & 7\,37084 \\
\hline
1058 & $32\,;\ 1, 8, 1, 3, 1, 3, (32)$; \ \ \ \ $34,\ 7, 1, 47,\ 14,\ 41,\ 17,\ (2)$ & 95685 \\
 & & 1318\,36323 \\
\hline
1059 & $32\,;\ 1, 31, 11, 25, 2, 1, 2, 3, 1, (31)$; \ \ \ \ $35,\ 1, 32, 5, 2, 23, 21, 35, 22, 15, 49,\ (2)$ & 53146 \\
 & & 68081 \\
\hline
1060 & $32\,;\ 1, 9, 5, 7, 3, 44, 1, 11, 1, 15, 9, (16)$; \ \ \ \ $36,\ 1, 28, 2, 1, 12, 3, 5, 1, 1, 2, 15, 1, 8, 2, 1, 2, (1, 1)$ & 38369 \\
 & & 737\,64856 \\
\hline
1061$\asterism$ & $32\,;\ 1, 9, 5, 7, 3, 44, 1, 11, 1, 15, 9, (16)$; \ \ \ \ $37,\ 27,\ 49,\ 44,\ 5,\ 20,\ 11,\ 35,\ 28,\ 25,\ 4,\ 55,\ 7,\ 23,\ 35,\ 20,\ (31, 31)$ & 82899 \\
 & & 11530 \\
\hline
1062 & $32\,;\ 1, 2, 3, 23, 19, (32)$; \ \ \ \ $38,\ (2)$ & 3\,06917 \\
 & & 9\,24122\,86971\,91521 \\
\hline
\end{longtable}

% ----- printed p. 446 (PDF 472) -----
\begin{center}
\small\textsc{Table 1001 to 1500---continued.}
\end{center}

\renewcommand{\arraystretch}{1.15}
\begin{longtable}{|r|p{0.55\textwidth}|r|}
\hline
$a$ & continued-fraction expansion of $\sqrt{a}$; period and subsidiary numbers & $x$ \\
 & & $y$ \\
\hline
\endhead
1063 & $32\,;\ 1, 1, 26, 25, 27, 37, 1, 10, 3, 1, 2, 1, 6, 1, 1, 21, 4, 1, (31)$; \ \ \ \ $39,\ 6,\ 17,\ 39,\ 18,\ 43,\ 9,\ 31,\ 34,\ 3,\ 13,\ 51,\ (2)$ & 12685 \\
 & & 5353\,15274 \\
\hline
1064 & $32\,;\ 1, 1, (28)$; \ \ \ \ $40,\ 24,\ 31,\ (1)$ & 21 \\
 & & 685 \\
\hline
1065 & $32\,;\ 1, 1, 23, 1, 2, 1, 3, 2, 1, 5, (4)$; \ \ \ \ $41,\ 23, 1, 35,\ 19,\ 39,\ 16,\ 21,\ 40,\ 11,\ (15)$ & 33160 \\
 & & 826\,67999 \\
\hline
1066$\asterism$ & $32\,;\ 1, 2, 2, 1, 5, 1, 6, (2, 2)$; \ \ \ \ $42,\ 39,\ 10,\ 49,\ 9,\ (25, 25)$ & 1\,34907 \\
 & & 05205 \\
\hline
1067 & $32\,;\ 1, (22)$; \ \ \ \ $43,\ 21, (1)$ & 3 \\
 & & 34 \\
\hline
1068 & $32\,;\ 1, 8, 49, 11, 5, (16)$; \ \ \ \ $44,\ 7, 1, (4)$ & 53094 \\
 & & 3 \\
\hline
1069$\asterism$ & $32\,;\ 1, 2, 3, 20, 17, 29, 31, 1, 21, 4, 3, 5, 7, 12, 1, 15, 2, 2, 1, 3, 1, 1, 1, 4, 1, (4, 4)$; \ \ \ \ $45,\ 6, 3, 15, 19, 12, 9, 5, 57, 4, 27, 20, 39, 15, 36, 25, 37, 12, 45, (13, 13)$ & 115\,39207 \\
 & & 186\,40986\,37841\,77726 \\
\hline
1070 & $32\,;\ 1, 2, 5, 26, 11, 31, 34, (5)$; \ \ \ \ $46, 1, 1, (12)$ & 6094\,77590\,16096\,85726\,12782 \\
 & & 98 \\
\hline
1071 & $32\,;\ 1, 18, 1, (6)$; \ \ \ \ $46, 17, 35, 1, 38, (9)$ & 90138 \\
 & & 48491 \\
\hline
1072 & $32\,;\ 1, 6, 1, 1, 8, 1, 4, 1, 1, (3)$; \ \ \ \ $47, 17, 35, 1, 38, (9)$ & 880 \\
 & & 28799 \\
\hline
1073$\asterism$ & $32\,;\ 1, (9, 9)$; \ \ \ \ $48, 16, 44, 9, 32, 33, 7, 49, 12, 33, 31, (16)$ & 1457\,20107 \\
 & & 81927 \\
\hline
1074 & $32\,;\ 1, 23, 31, 30, 25, (2)$; \ \ \ \ $49, 15, 3, (7, 7)$ & 47710\,1385 \\
 & & 45368 \\
\hline
1075 & $32\,;\ 50, 3, 1, 2, 3, 10, 1, 1, 1, 2, 2, 6, (1)$; \ \ \ \ $51, 14, 1, 19, 6, 39, 25, 34, 21, 26, 9, (50)$ & 1\,06476 \\
 & & 12729 \\
\hline
1076 & $32\,;\ 51, 14, 1, (3, 39, 21)$; \ \ \ \ $52, 4, 2, 15, (16)$ & 34\,89425 \\
 & & 16\,51504 \\
\hline
1077 & $32\,;\ 52, 4, 2, 15, (16)$; \ \ \ \ $53, 29, 4, 1, (20)$ & 88675 \\
 & & 410 \\
\hline
1078 & $32\,;\ 53, 29, 4, 1, (20)$; \ \ \ \ $54, 4, 10, 1, 1, 2, 2, 4, 3, 1, (1)$ & 74226 \\
 & & 13449 \\
\hline
1079 & $32\,;\ 54, 4, 10, 1, 1, 2, 2, 4, 3, 1, (1)$; \ \ \ \ $55, 1, 35, 29, (3)$ & 22399 \\
 & & 2357\,16760 \\
\hline
1080 & $32\,;\ 55, 1, 35, 29, (3)$; \ \ \ \ $56, 9, 5, (6)$ & 1805 \\
 & & 76876 \\
\hline
1081 & $32\,;\ 7, 1, 4, 3, 1, 6, 1, 1, 5, 1, 3, 1, 12, 2, 1, 1, 1, 21, 3, 2, (2)$; \ \ \ \ $57,\ 8,\ 15,\ 16,\ 45,\ 9,\ 33,\ 32,\ 11,\ 27,\ 15,\ 48,\ 5,\ 24,\ 33,\ 25,\ 40,\ 3,\ 19,\ 24,\ (23)$ & 97325 \\
 & & 54\,161 \\
\hline
1082$\asterism$ & $32\,;\ (2, 7, 1, 5)$; \ \ \ \ $58,\ 7,\ 26,\ 23,\ 22,\ (31, 31)$ & 8918\,12221\,87280 \\
 & & 2\,93215\,05610\,17151\,07648 \\
\hline
1083 & $32\,;\ (9)$; \ \ \ \ $59,\ (6)$ & 19615 \\
 & & 12\,38369 \\
\hline
\end{longtable}

% ----- printed p. 447 (PDF 473) -----
\begin{center}
\small\textsc{Table 1001 to 1500---continued.}
\end{center}

\renewcommand{\arraystretch}{1.15}
\begin{longtable}{|r|p{0.55\textwidth}|r|}
\hline
$a$ & continued-fraction expansion of $\sqrt{a}$; period and subsidiary numbers & $x$ \\
 & & $y$ \\
\hline
\endhead
1084 & $32\,;\ 1, 1, 1, 3, 1, 2, 1, 1, 6, 1, 1, 1, 12, 5, 1, 1, 1, 21, 3, 4, (16)$; \ \ \ \ $60,\ 5,\ 12,\ 25,\ 24,\ 17,\ 27,\ 9,\ 51,\ 8,\ 23,\ 36,\ 19,\ 40,\ 15,\ 37,\ 24,\ 41,\ 3,\ 20,\ 15,\ (4)$ & 8\,17041\,27029\,09269\,93200 \\
 & & 269\,00393\,64201\,05379\,19999 \\
\hline
1085 & $32\,;\ 1, 15, (2)$; \ \ \ \ $61,\ 4,\ (31)$ & 544 \\
 & & 17919 \\
\hline
1086 & $32\,;\ 1, (20)$; \ \ \ \ $62,\ (3)$ & 22 \\
 & & 725 \\
\hline
1087 & $32\,;\ 1, (31)$; \ \ \ \ $63,\ (2)$ & 33 \\
 & & 1088 \\
\hline
1088 & $32\,;\ (1)$; \ \ \ \ $(64)$ & 1 \\
 & & 33 \\
\hline
1090$\asterism$ & $33\,;\ (66)$; \ \ \ \ $(1)$ & 33 \\
 & & 1090 \\
\hline
1091 & $33\,;\ (33)$; \ \ \ \ $(2)$ & 33 \\
 & & 2 \\
\hline
1092 & $33\,;\ (22)$; \ \ \ \ $(3)$ & 727 \\
 & & 22 \\
\hline
1093$\asterism$ & $33\,;\ 16, (1, 1)$; \ \ \ \ $4,\ (33, 33)$ & 18018 \\
 & & 545 \\
\hline
1094 & $33\,;\ 13, 4, 1, 1, 1, 5, 2, 1, 2, 3, 9, 6, 1, 1, (32)$; \ \ \ \ $5,\ 14,\ 37,\ 25,\ 38,\ 11,\ 23,\ 35,\ 22,\ 19,\ 7,\ 10,\ 31,\ 35,\ (2)$ & 45043\,80474\,67914 \\
 & & 14\,89854\,05815\,38085 \\
\hline
1095 & $33\,;\ (11)$; \ \ \ \ $(6)$ & 364 \\
 & & 11 \\
\hline
1096 & $33\,;\ 9, 2, 3, 1, (15)$; \ \ \ \ $7,\ 28,\ 15,\ 49,\ (4)$ & 1\,19595 \\
 & & 39\,59299 \\
\hline
1097$\asterism$ & $33\,;\ 8, 3, 1, 3, 2, 1, (1, 1)$; \ \ \ \ $8,\ 17,\ 41,\ 16,\ 23,\ 32,\ (29, 29)$ & 2106\,34621 \\
 & & 19276 \\
\hline
1098 & $33\,;\ 7, 2, 1, (6)$; \ \ \ \ $9,\ 22,\ 41,\ (9)$ & 3564 \\
 & & 1\,18097 \\
\hline
1099 & $33\,;\ 6, 1, 1, 1, 1, 2, 21, 1, 2, 1, 1, 6, 1, (3)$; \ \ \ \ $10,\ 37,\ 27,\ 30,\ 31,\ 25,\ 3,\ 46,\ 19,\ 30,\ 35,\ 9,\ 47,\ (14)$ & 4\,80575\,45715 \\
 & & 199 \\
\hline
1100 & $33\,;\ (6)$; \ \ \ \ $(11)$ & 6 \\
 & & 199 \\
\hline
1101 & $33\,;\ 5, 1, 1, 16, (22)$; \ \ \ \ $12,\ 31,\ 35,\ 4,\ (3)$ & 7\,32732 \\
 & & 243\,13015 \\
\hline
1102 & $33\,;\ 5, 10, 1, 6, (2)$; \ \ \ \ $13,\ 6,\ 53,\ 9,\ (29)$ & 3\,42882 \\
 & & 113\,82443 \\
\hline
1103 & $33\,;\ 4, 1, 2, 1, 2, 3, 1, 1, 5, 2, 9, (33)$; \ \ \ \ $14,\ 41,\ 19,\ 37,\ 22,\ 17,\ 31,\ 34,\ 11,\ 29,\ 7,\ (2)$ & 7\,06456\,11145 \\
 & & 234\,62427\,45024 \\
\hline
1104 & $33\,;\ 4, 2, (2)$; \ \ \ \ $15,\ 25,\ (23)$ & 234 \\
 & & 7775 \\
\hline
1105$\asterism$ & $33\,;\ 4, (7, 7)$; \ \ \ \ $16,\ (9, 9)$ & 857 \\
 & & 28488 \\
\hline
\end{longtable}

% ----- printed p. 448 (PDF 474) -----
\begin{center}
\small\textsc{Table 1001 to 1500---continued.}
\end{center}

\renewcommand{\arraystretch}{1.15}
\begin{longtable}{|r|p{0.55\textwidth}|r|}
\hline
$a$ & continued-fraction expansion of $\sqrt{a}$; period and subsidiary numbers & $x$ \\
 & & $y$ \\
\hline
\endhead
1106 & $33\,;\ 1, 1, (7)$; \ \ \ \ $17,\ 7,\ 2,\ (8)$ & 152 \\
 & & 5055 \\
\hline
1107 & $33\,;\ 3, 46, 23, 9, (33)$; \ \ \ \ $18,\ 3,\ 37,\ (16)$ & 18295 \\
 & & 19428 \\
\hline
1108 & $33\,;\ 3, 49, 16, 13, 48, 9, 36, 29, 23, 21, 39, 16, 27, 12, 7, 21, (4)$; \ \ \ \ $19,\ 28,\ 3,\ 4,\ 1,\ 6,\ 1,\ 1,\ 2,\ 2,\ 1,\ 3,\ 2,\ 5,\ 9,\ 3,\ (16)$ & 63026 \\
 & & 4781\,20058\,69390\,13510 \\
\hline
1109$\asterism$ & $33\,;\ 3, 11, 44, 3, 13, 16, 1, 1, 2, 1, 4, (2, 2)$; \ \ \ \ $20,\ 3,\ 19,\ 5,\ 1,\ 17,\ 5,\ 4,\ 37,\ 29,\ 20,\ 41,\ 13,\ (25, 25)$ & 35957\,38617 \\
 & & 1832\,42610 \\
\hline
1110 & $33\,;\ 3, (6)$; \ \ \ \ $21,\ (10)$ & 61020\,60015 \\
 & & 60 \\
\hline
1111 & $33\,;\ (3)$; \ \ \ \ $(22)$ & 1003 \\
 & & 30 \\
\hline
1112 & $33\,;\ 2, (7)$; \ \ \ \ $23,\ (8)$ & 21056 \\
 & & 2501 \\
\hline
1113 & $33\,;\ 2, 17, 16, 47, (7)$; \ \ \ \ $24,\ 41,\ 2,\ 3,\ 3,\ 1,\ (8)$ & 75 \\
 & & 7\,02463 \\
\hline
1114$\asterism$ & $33\,;\ 7, 55, 6, 15, 22, 39, 15, 42, 17, 9, 10, (33, 33)$; \ \ \ \ $25,\ 37,\ 2,\ 39,\ 8,\ 1,\ 10,\ 4,\ 2,\ 1,\ 3,\ 1,\ 3,\ 7,\ 6,\ (1, 1)$ & 18311\,57471\,56745 \\
 & & 6\,11178\,81032\,02293 \\
\hline
1115 & $33\,;\ 5, 2, 2, 2, 5, 1, (5)$; \ \ \ \ $26,\ 29,\ 3,\ 1,\ 4,\ (5)$ & 45 \\
 & & 1\,36565 \\
\hline
1116 & $33\,;\ 5, 1, 11, 40, 23, (4)$; \ \ \ \ $27,\ 25,\ 35,\ 11,\ 11,\ 49,\ (10)$ & 46\,60126 \\
 & & 1\,20799 \\
\hline
1117$\asterism$ & $33\,;\ 5, 1, 1, 7, 5, 4, 1, 21, 2, 9, 16, 1, 1, 1, 1, 6, 1, 4, 1, (2, 2)$; \ \ \ \ $28,\ 1,\ 2,\ 1,\ 5,\ 36,\ 23,\ 12,\ 13,\ 52,\ 3,\ 31,\ 7,\ 4,\ 39,\ 27,\ 28,\ 37,\ 9,\ 49,\ 12,\ 41,\ (21, 21)$ & 38320 \\
 & & 6272\,59559\,53855\,43645 \\
\hline
1118 & $33\,;\ 5, (2)$; \ \ \ \ $29,\ (26)$ & 126 \\
 & & 4213 \\
\hline
1119 & $33\,;\ 5, 38, 25, 39, 10, (3)$; \ \ \ \ $30,\ 4,\ 1,\ 1,\ 1,\ 6,\ (22)$ & 9 \\
 & & 316\,46364 \\
\hline
1120 & $33\,;\ 5, (4)$; \ \ \ \ $31,\ 1, (15)$ & 57255 \\
 & & 1 \\
\hline
1121 & $33\,;\ 9, 1, 1, 8, 1, 1, 1, 3, 1, 1, 7, 1, 4, 3, 1, (2)$; \ \ \ \ $32,\ 56,\ 7,\ 40,\ 25,\ 37,\ 16,\ 31,\ 35,\ 8,\ 49,\ 13,\ 17,\ 40,\ (19)$ & 453\,3765 \\
 & & 15198\,93975\,58620 \\
\hline
1122 & $33\,;\ (1)$; \ \ \ \ $(33)$ & 67 \\
 & & 51043\,64951 \\
\hline
1123 & $33\,;\ 4, 1, 21, 1, 5, 7, 3, 1, 1, 2, 2, 10, 1, 3, (33)$; \ \ \ \ $34,\ 33,\ 3,\ 54,\ 11,\ 9,\ 18,\ 33,\ 31,\ 22,\ 27,\ 6,\ 47,\ 17,\ (2)$ & 49257\,14799 \\
 & & 16\,50664\,55626\,32482 \\
\hline
1124 & $33\,;\ 9, 13, 3, 3, 1, 1, 1, 1, 1, 2, (16)$; \ \ \ \ $35,\ 7,\ 5,\ 20,\ 17,\ 35,\ 28,\ 31,\ 29,\ 32,\ 25,\ (4)$ & 2266\,74757 \\
 & & 22006\,30049 \\
\hline
1125 & $33\,;\ 11, 36, 5, 1, 29, 25, 9, 4, 41, 25, 36, 19, 44, (9)$; \ \ \ \ $36,\ 31,\ 2718,\ 91205,\ 173\,04001$ & 5\,16002 \\
 & & 173\,04001 \\
\hline
1126 & $33\,;\ 3, 1, (32)$; \ \ \ \ $37,\ 30,\ (2)$ & 2718 \\
 & & 91205 \\
\hline
\end{longtable}

% ----- printed p. 449 (PDF 475) -----
\begin{center}
\small\textsc{Table 1001 to 1500---continued.}
\end{center}

\renewcommand{\arraystretch}{1.15}
\begin{longtable}{|r|p{0.55\textwidth}|r|}
\hline
$a$ & continued-fraction expansion of $\sqrt{a}$; period and subsidiary numbers & $x$ \\
 & & $y$ \\
\hline
\endhead
1127 & $33\,;\ 1, 3, (33)$; \ \ \ \ $38,\ 29,\ 19,\ (2)$ & 1645 \\
 & & 55224 \\
\hline
1128 & $33\,;\ 1, 1, 4, (3)$; \ \ \ \ $39,\ 28,\ 13,\ (24)$ & 70 \\
 & & 2351 \\
\hline
1129$\asterism$ & $33\,;\ 1, 1, (1)$; \ \ \ \ $40,\ 27,\ (37)$ & 5 \\
 & & 168 \\
\hline
1130$\asterism$ & $33\,;\ 1, 1, (1, 1)$; \ \ \ \ $41,\ 25,\ (31, 31)$ & 13 \\
 & & 437 \\
\hline
1131 & $33\,;\ 1, 1, 5, 2, 1, 1, 1, (4)$; \ \ \ \ $42,\ 13,\ 23,\ 35,\ 30,\ 35,\ (13)$ & 10948 \\
 & & 3\,68185 \\
\hline
1132 & $33\,;\ 1, 1, 1, 4, 1, 16, 3, 1, 1, 21, 1, 7, 2, 5, 7, 3, 2, 2, (16)$; \ \ \ \ $43,\ 24,\ 1,\ 23,\ 31,\ 36,\ 3,\ 57,\ 8,\ 29,\ 12,\ 9,\ 19,\ 24,\ 27,\ (4)$ & 1\,13654\,94862\,36362 \\
 & & 38\,23944\,35058\,85447 \\
\hline
1133 & $33\,;\ 1, 23, 39, 13, 43, (6)$; \ \ \ \ $44,\ 2,\ (11)$ & 107\,14999 \\
 & & 15300 \\
\hline
1134 & $33\,;\ 1, 22, 4, 7, 2, 2, (4)$; \ \ \ \ $45,\ 21,\ 25,\ 9,\ 29,\ 26,\ 25,\ (14)$ & 52041 \\
 & & 3620\,74049 \\
\hline
1135 & $33\,;\ 1, 4, 10, 1, 1, 3, 1, 1, 2, 1, 1, 1, 6, 1, (5)$; \ \ \ \ $46,\ 20,\ 1,\ 23,\ 6,\ 49,\ 15,\ 34,\ 31,\ 21,\ 35,\ 26,\ 39,\ 9,\ 51,\ (10)$ & 31\,50736\,19505 \\
 & & 1061\,47549 \\
\hline
1136 & $33\,;\ 1, 14, 1, 2, 9, (4)$; \ \ \ \ $47,\ 19,\ 32,\ 31,\ 25,\ 7,\ (16)$ & 56124 \\
 & & 2427\,18620 \\
\hline
1137 & $33\,;\ 1, 3, 2, 1, 1, 7, 1, 5, 4, (22)$; \ \ \ \ $48,\ 18,\ 17,\ 17,\ 24,\ 29,\ 37,\ 8,\ 51,\ 11,\ 16,\ (3)$ & 20799 \\
 & & 26292 \\
\hline
1138$\asterism$ & $33\,;\ 4, 1, (1, 1)$; \ \ \ \ $49,\ 41,\ 14,\ (33, 33)$ & 10337 \\
 & & 3\,48711 \\
\hline
1139 & $33\,;\ (1)$; \ \ \ \ $(17)$ & 1354 \\
 & & 50\,16,\ (15) \\
\hline
1140 & $33\,;\ (15)$; \ \ \ \ $(4)$ & 2431 \\
 & & 51 \\
\hline
1141 & $33\,;\ 3, 1, 1, 15, 31, 36, 1, 12, 1, 21, 1, 1, 2, 1, 4, 3, 7, 5, 16, 1, 2, 3, 1, 1, 1, 2, 1, 4, 1, (8)$; \ \ \ \ $52,\ 5,\ 60,\ 3,\ 39,\ 28,\ 25,\ 12,\ 15,\ 20,\ 9,\ 13,\ 4,\ 45,\ 21,\ 17,\ 36,\ 27,\ 35,\ 20,\ 39,\ 19,\ 1,\ 43,\ 12,\ 51,\ (7)$ & 3\,06933\,85322\,76565\,71973 \\
 & & 103\,67823\,94157\,22396\,32371\,97208 \\
\hline
1142 & $33\,;\ 1, 14, 11, 2, 1, 8, 1, (32)$; \ \ \ \ $53,\ 13,\ 47,\ 1,\ 22,\ 43,\ 7,\ 59,\ (2)$ & 268\,25215 \\
 & & 9066\,12101 \\
\hline
1143 & $33\,;\ 1, 3, 14, 5, 1, 1, 2, (3)$; \ \ \ \ $54,\ 27,\ 34,\ 23,\ (18)$ & 47\,28010 \\
 & & 72 \\
\hline
1144 & $33\,;\ 12, 4, 4, 25, 1, 37, 1, 6, (1)$; \ \ \ \ $55,\ 39,\ 40,\ 9,\ (52)$ & 1\,39925 \\
 & & 37 \\
\hline
1145$\asterism$ & $33\,;\ 4, (5, 5)$; \ \ \ \ $56,\ (15, 15)$ & 30624 \\
 & & 1252 \\
\hline
1146 & $33\,;\ 10, 47, 5, 1, 1, 1, (10)$; \ \ \ \ $57,\ 1,\ 38,\ 25,\ 41,\ (6)$ & 16611 \\
 & & 5\,61835 \\
\hline
\end{longtable}

\end{document}
